\documentclass[11pt,a4paper]{article}
\usepackage{amsmath,amssymb}
\title{On Sorting Large Collections}
\author{Priya Nandakumar}
\date{September 9, 2022}
\begin{document}
\maketitle

\begin{abstract}
We review the classical comparison-based sorting algorithms and discuss the
information-theoretic lower bound that constrains all of them.
\end{abstract}

\section{Introduction}
Sorting is among the most thoroughly studied problems in computer science
\cite{knuth}. Textbook treatments such as \cite{clrs} describe dozens of
methods, yet only a handful are used in practice. Detailed analysis of
quicksort appears in \cite[p.~5]{knuth}, and a historical account of the
subject can be found by consulting \cite{knuth,clrs} together.

\section{A Lower Bound}
Any comparison sort must distinguish among $n!$ possible orderings. Each
comparison yields one bit, so at least $\log_2 (n!)$ comparisons are needed in
the worst case. By Stirling's approximation,
\begin{equation}\label{eq:stirling}
\log_2 (n!) = n\log_2 n - n\log_2 e + O(\log n).
\end{equation}
Hence the number of comparisons is $\Omega(n \log n)$, as claimed in
Section~\ref{sec:practice}.

\section{Practical Considerations}\label{sec:practice}
Real implementations care about caching and memory movement as much as about
comparison counts. A simple verbatim sketch of insertion sort follows.
\begin{verbatim}
for i = 1 to n-1:
    key = a[i]
    j = i - 1
    while j >= 0 and a[j] > key:
        a[j+1] = a[j]
        j = j - 1
    a[j+1] = key
\end{verbatim}
The inner loop uses the comparison \verb|a[j] > key|, which dominates the
running time for nearly reversed input.

\begin{thebibliography}{99}
\bibitem[Knu84]{knuth} Donald E. Knuth, \emph{The Art of Computer
  Programming, Volume 3: Sorting and Searching}, Addison-Wesley, 1984.
\bibitem{clrs} Thomas H. Cormen, Charles E. Leiserson, Ronald L. Rivest and
  Clifford Stein, \emph{Introduction to Algorithms}, MIT Press, 2009.
\end{thebibliography}
\end{document}
